\chapter{Trade Capture and Booking}\label{ch:pl:trade-capture-and-booking}

A swap of \$100~million was amended to \$150~million on Monday, novated to another dealer on Wednesday and partially terminated on Friday
morning. On Friday afternoon the risk system held version~3 (\$150~million, the new dealer), the confirmation system version~2 (\$150~million,
the old dealer) and the settlement system version~4 (\$90~million, the new dealer). Each had copied the trade when something it cared about
happened, and none of them had followed what happened to it. This chapter treats a trade as the result of its events: captured once in a
canonical model, booked by rule, allocated by a stated rule, and changed only by events that every downstream system reads.

\section{What a trade is, as data}

\begin{definition}[Canonical trade model]\label{def:pl:trade-capture-and-booking:canonical}
A \emph{canonical trade model}\index{canonical trade model} is the firm's single representation of a trade shared by every system that
handles it: a header (identifiers, parties, book, dates, status, version) and the economics (the instrument, quantity and price), with
rules that every valid trade satisfies.
\end{definition}

\begin{definition}[Unique transaction identifier]\label{def:pl:trade-capture-and-booking:uti}
A \emph{unique transaction identifier}\index{unique transaction identifier} (UTI) identifies one transaction across all the parties and
repositories that report it: under the CPMI--IOSCO technical guidance, the legal entity identifier (Book~2, chapter~28) of the entity that
generates it followed by a value unique for that entity, at most 52 characters, upper-case letters and digits only.
\end{definition}

The chapter's model (\texttt{firm.tradecap}) keeps the instrument in the form Book~5's pricing library serialises it, so that a captured trade
is priced, risk-managed and reported from the same record; the validator checks the identifiers' formats, that buyer and seller differ, that
settlement does not precede the trade, and that the quantity is positive. The industry's own models --- FpML for messages, ISDA's Common Domain
Model for events --- follow the same idea at a much larger scale.

\section{Capture: from execution to booked trade}

\begin{definition}[Trade capture, booking]\label{def:pl:trade-capture-and-booking:capture}
\emph{Trade capture}\index{trade capture} turns an execution --- an exchange execution report, an OTC ticket, a voice trade --- into a trade
in the canonical model, once and as early as possible. \emph{Booking}\index{booking} assigns the captured trade to a book, the unit that owns
its risk and P\&L.
\end{definition}

\begin{definition}[Booking model]\label{def:pl:trade-capture-and-booking:model}
A \emph{booking model}\index{booking model} is the set of rules that decides where a trade is booked --- by instrument type, desk, legal
entity and strategy --- so that the book is a consequence of the trade's attributes, not a field typed by hand.
\end{definition}

The chapter's exchange side is a block: 9\,900 shares bought on Book~10's simulator by an execution algorithm sending an order that takes the
offer every 20~seconds, filled in 57 executions over three minutes at an average of \$100.0783. Each execution report becomes a trade with its
own UTI, bought by the firm from the clearing house. The OTC side is 200 swaps from six dealers, captured from tickets on Monday morning.

\begin{omfigure}
\begin{tikzpicture}
\begin{axis}[width=11.5cm,height=5cm,xlabel={seconds after the open},ylabel={price (\$)},grid=major,grid style={gray!20},
  tick label style={font=\small},label style={font=\small},table/col sep=comma,
  scatter,only marks,point meta={\thisrow{qty}},scatter/use mapped color={draw=omDef,fill=omDef!40},
  visualization depends on={sqrt(\thisrow{qty})/6 \as \msize},scatter/@pre marker code/.append style={/tikz/mark size=\msize}]
\addplot table[x=seconds,y=price]{figdata/platforms/21-trade-capture-and-booking/block.csv};
\draw[omThm,dashed] (axis cs:0,100.0783) -- (axis cs:220,100.0783) node[pos=0.02,above right,font=\scriptsize]{average \$100.0783};
\end{axis}
\end{tikzpicture}

\omcaption{The 57 executions of the 9\,900-share block on the simulator, by time and price (marker area by size): ten child orders taking the
offer every 20~seconds, each filled in several executions as it walks the book. Every execution is captured as a trade; every fund is allocated
at the average. Data: \texttt{fig\_tradecap.py}.}\label{fig:pl:trade-capture-and-booking:block}
\end{omfigure}

\section{The lifecycle as events}

\begin{definition}[Trade lifecycle event, trade amendment]\label{def:pl:trade-capture-and-booking:event}
A \emph{trade lifecycle event}\index{trade lifecycle event} is anything that happens to a trade after it is executed: an amendment, a
novation (Book~1, chapter~5), a partial termination, an exercise, an allocation, a cancellation. A \emph{trade amendment}\index{trade amendment}
changes a term of the trade (notional, rate, dates) by agreement; it is recorded as a new version, never by overwriting the old one.
\end{definition}

ISDA's Common Domain Model describes lifecycle events in the same terms: ``a lifecycle event describes a state transition,'' with a
before and an after state, composed from a few primitive instructions --- execution, quantity change, terms change, party change, exercise,
transfer and a handful of others. \texttt{firm.tradecap} keeps the events and computes the trade: a store applies events in time order and
rebuilds any version as of any time (\cref{lst:pl:trade-capture-and-booking:store}). The event sourcing of chapter~24 is the general
pattern; here it is the only way to answer, on Friday, what the trade was on Wednesday.

\omcode{code/firm/tradecap/firm_tradecap.py}{70}{98}{The trade store: events appended, never edited; a trade as of any time is the replay
of its events up to then.}\label{lst:pl:trade-capture-and-booking:store}

\omcode{code/firm/tradecap/firm_tradecap.py}{101}{124}{Each event kind as a transition from one version of the trade to the
next.}\label{lst:pl:trade-capture-and-booking:step}

\begin{omfigure}
\begin{tikzpicture}[font=\scriptsize,>=Stealth,ver/.style={draw,rounded corners=2pt,align=center,minimum height=1.0cm,minimum width=2.2cm}]
\node[ver,fill=gray!10] (v1) at (0,0) {v1: \$100m\\ dealer 1};
\node[ver,fill=gray!10] (v2) at (3.4,0) {v2: \$150m\\ dealer 1};
\node[ver,fill=gray!10] (v3) at (6.8,0) {v3: \$150m\\ dealer 4};
\node[ver,fill=omDef!15] (v4) at (10.2,0) {v4: \$90m\\ dealer 4};
\draw[->] (v1) -- node[above=4pt,align=center]{amend\\ Mon} (v2);
\draw[->] (v2) -- node[above=4pt,align=center]{novate\\ Wed} (v3);
\draw[->] (v3) -- node[above=4pt,align=center]{cut 40\%\\ Fri} (v4);
\node[omThm,align=center] at (3.4,-1.3) {confirmation\\ (copies on new, amend)};
\node[omThm,align=center] at (6.8,-1.3) {risk\\ (copies nightly)};
\node[omThm,align=center] at (10.2,-1.3) {settlement\\ (copies on new, terminate)};
\draw[->,omThm] (3.4,-0.95) -- (v2.south); \draw[->,omThm] (6.8,-0.95) -- (v3.south); \draw[->,omThm] (10.2,-0.95) -- (v4.south);
\end{tikzpicture}

\omcaption{The hook's swap through the week and what three copy-based systems held on Friday at 15:00. Each copied the trade on the events
it cares about, or on its own schedule. Settlement holds version~4 only because the last event, a termination, is one it copies on; only a
system that follows every event holds the current version whatever the last event was.}\label{fig:pl:trade-capture-and-booking:versions}
\end{omfigure}

\section{Three versions on Friday}

The week's lifecycle is seeded: 69 amendments, 26 novations and 41 partial terminations among the 200 swaps between Monday 10:00 and Friday
14:00, the hook's swap among them. Three downstream systems keep copies: risk copies every trade at 18:00 each evening; confirmation receives a copy on
the events it confirms (new trades and amendments); settlement on the events that move cash (new trades and terminations). On Friday at 15:00
their copies are compared with the event-sourced trade.

\omcode{code/platforms/21-trade-capture-and-booking/python/pl_tradecap.py}{118}{144}{The three copy policies and the comparison with the
event-sourced trade on Friday afternoon.}\label{lst:pl:trade-capture-and-booking:downstream}

\begin{table}[ht]
\centering\small
\begin{tabular}{lrr}
\toprule
system & wrong, copying & wrong, following events \\
\midrule
risk (nightly copy) & 25 & 0 \\
confirmation (copy on new and amend) & 60 & 0 \\
settlement (copy on new and terminate) & 71 & 0 \\
\midrule
trades on which the three disagree & 109 of 200 & 0 \\
\bottomrule
\end{tabular}
\caption{Swaps whose downstream copy differs from the event-sourced trade on Friday at 15:00, after a week of lifecycle
events.}\label{tab:pl:trade-capture-and-booking:downstream}
\end{table}

\begin{omfigure}
\begin{tikzpicture}
\begin{axis}[width=11.5cm,height=5cm,xbar,bar width=8pt,xmin=0,xmax=130,xlabel={swaps wrong on Friday at 15:00 (of 200)},
  symbolic y coords={any disagreement,settlement,confirmation,risk},ytick=data,enlarge y limits=0.2,
  tick label style={font=\small},label style={font=\small},grid=major,grid style={gray!20},table/col sep=comma,
  nodes near coords,nodes near coords style={font=\scriptsize,anchor=west},legend style={font=\footnotesize,at={(0.5,1.03)},
  anchor=south,legend columns=2}]
\addplot[fill=omThm!45,draw=omThm] table[y=system,x=copy_wrong]{figdata/platforms/21-trade-capture-and-booking/downstream.csv};
\addlegendentry{copying}
\addplot[fill=omDef!45,draw=omDef] table[y=system,x=events_wrong]{figdata/platforms/21-trade-capture-and-booking/downstream.csv};
\addlegendentry{following the events}
\end{axis}
\end{tikzpicture}

\omcaption{Downstream systems against the event-sourced swaps on Friday afternoon, after a week of amendments, novations and partial
terminations: copies are wrong in 25 to 71 trades and disagree on 109; systems that follow the events are never wrong. Data:
\texttt{fig\_tradecap.py}.}\label{fig:pl:trade-capture-and-booking:downstream}
\end{omfigure}

Copying fails in a different way in each system (\cref{tab:pl:trade-capture-and-booking:downstream}): risk is wrong only about the day's events
(25 trades), confirmation about every novation and termination it was never sent (60), settlement about every amendment and novation (71), and
on 109 of the 200 swaps at least two systems disagree. Following the events, all three hold the same trade, the one the store rebuilds.

\section{Allocations and straight-through processing}

\begin{definition}[Block allocation]\label{def:pl:trade-capture-and-booking:allocation}
A \emph{block allocation}\index{block allocation} splits a trade executed for several accounts at once --- a block for several funds --- among
them after execution, at the block's average price, by a rule fixed before the trade that also decides where the units that do not divide
evenly go.
\end{definition}

\omcode{code/firm/tradecap/firm_tradecap.py}{144}{166}{Allocation in whole lots at the average price, with three rules for the lots that do
not divide.}\label{lst:pl:trade-capture-and-booking:allocate}

The block of 9\,900 shares (99 lots of 100) is allocated to three funds with targets of 50\%, 30\% and 20\%, whose exact shares are 4\,950, 2\,970
and 1\,980. The largest-remainder rule gives 4\,900, 3\,000 and 2\,000: all 99 lots allocated, no fund more than 50 shares from its target. Giving
the odd lots to the largest funds gives 5\,000, 3\,000 and 1\,900, 80 shares from target for the smallest. Rounding each fund on its own gives
5\,000, 3\,000 and 2\,000 --- 10\,000 shares for a block of 9\,900: an over-allocation of one lot that someone must buy or take back.
(\cref{tab:pl:trade-capture-and-booking:allocation}). Every fund pays the block's average price, \$100.0783, whatever executions made it.

\begin{table}[ht]
\centering\small
\begin{tabular}{lrrrrr}
\toprule
rule & fund A & fund B & fund C & residual & largest deviation \\
\midrule
largest remainder & 4\,900 & 3\,000 & 2\,000 & 0 & 50 \\
odd lots to the largest funds & 5\,000 & 3\,000 & 1\,900 & 0 & 80 \\
each fund rounded & 5\,000 & 3\,000 & 2\,000 & $-100$ & 50 \\
\bottomrule
\end{tabular}
\caption{A block of 9\,900 shares allocated to three funds (targets 50\%, 30\%, 20\%) in lots of 100 under three rules: shares per fund, the
block minus the allocation, and the largest deviation from a fund's exact share.}\label{tab:pl:trade-capture-and-booking:allocation}
\end{table}

\begin{definition}[Straight-through processing]\label{def:pl:trade-capture-and-booking:stp}
\emph{Straight-through processing}\index{straight-through processing} carries a trade from execution to settlement --- capture, \omterm{def:pl:trade-capture-and-booking:capture}{booking},
allocation, confirmation, clearing, settlement --- without manual re-entry or intervention, each step reading the output of the previous one in
the canonical model.
\end{definition}

\omterm{def:pl:trade-capture-and-booking:stp}{Straight-through processing} is what the canonical model and the event stream buy: each manual step is a copy, and each copy is a version that can
drift. The rate of trades that pass through without a human touch is the operational measure of a platform's \omterm{def:pl:trade-capture-and-booking:capture}{trade capture} (chapter~22 measures
the breaks when they do not).

\begin{dated}{2026-09}{Identifiers and event models}\label{dat:pl:trade-capture-and-booking:standards}
CPMI and IOSCO published their technical guidance on the harmonisation of the UTI on 28 February 2017: the generating entity's LEI concatenated
with a value unique for that entity, no check digit, at most 52 characters from A--Z and 0--9. ISDA's Common Domain Model, maintained as open
source, models lifecycle events as state transitions composed from primitive instructions (execution, quantity change, terms change, party
change, exercise and others).
\end{dated}

\section{Tutorial: one block, two hundred swaps, one week}\label{tut:pl:trade-capture-and-booking:tut}

\begin{tutorial}
\textbf{Goal.} Capture executions and tickets into one model, allocate a block, run a week of lifecycle events, and compare copying with
following events. \textbf{End state:} \cref{tab:pl:trade-capture-and-booking:downstream,tab:pl:trade-capture-and-booking:allocation}.
\begin{tutsteps}
\item \textbf{Capture}: \texttt{pl\_tradecap.block\_fills()} and \texttt{capture\_block}; validate every trade.
\item \textbf{Allocate}: \texttt{allocations(fills)} under the three rules.
\item \textbf{Lifecycle}: \texttt{week()}; \texttt{store.versions} and \texttt{store.as\_of}.
\item \textbf{Downstream}: \texttt{downstream(store, events)}.
\end{tutsteps}
\textbf{What to change next.} Add an exercise event that turns a swaption into a swap and follow it downstream; allocate by the funds' cash
instead of fixed targets.
\end{tutorial}

\section{Build: trade capture}\label{bld:pl:trade-capture-and-booking:tradecap}

\begin{build}
\bfield{Purpose} One trade record for every system, changed only by events, from capture to settlement.
\bfield{Interface} \texttt{uti}, \texttt{validate}, \texttt{Event}, \texttt{TradeStore} (\texttt{apply}, \texttt{as\_of},
\texttt{versions}), \texttt{from\_execution}, \texttt{from\_ticket}, \texttt{allocate}, \texttt{BookingModel}.
\bfield{Rules} Trades are never edited, only evented; every trade validates; identifiers follow the UTI and LEI formats; allocation rules are
fixed before the trade and allocate the whole block; books come from rules.
\bfield{Acceptance tests} \texttt{code/firm/tradecap/tests/}: UTI format and validation errors; a lifecycle's versions and states as of any
time, including termination to zero and an event on an unknown trade; exercise and capture from an execution report; the three allocation
rules; the \omterm{def:pl:trade-capture-and-booking:model}{booking model}.
\bfield{Stretch} Messages in FpML or events in the CDM's JSON; allocations by cash with fractional shares where allowed; a straight-through
rate report.
\end{build}

\begin{omsources}
\item CPMI and IOSCO, \emph{Harmonisation of the Unique Transaction Identifier: Technical Guidance}, February 2017.
\item ISDA Common Domain Model documentation, ``Event Model''.
\item One Quant Book~1, chapter~5, and Book~2, chapter~28 (novation, LEIs); One Quant Book~10, chapter~26 (execution reports).
\end{omsources}

\section{Exercises}

\begin{exercise}[$\star$]\label{exo:pl:trade-capture-and-booking:1}
A block of 1\,000 shares is split among three funds with equal targets in lots of 100. What does each rule give?
\end{exercise}

\begin{exercise}[$\star$]\label{exo:pl:trade-capture-and-booking:2}
Why is the risk system wrong about fewer trades than the other two?
\end{exercise}

\begin{exercise}[$\star$]\label{exo:pl:trade-capture-and-booking:3}
What does a UTI generated by the firm for its 7th trade look like?
\end{exercise}

\begin{exercise}[$\star\star$]\label{exo:pl:trade-capture-and-booking:4}
Why must the allocation rule be fixed before the trade, and what goes wrong if a trader chooses it after seeing the fills?
\end{exercise}

\begin{exercise}[$\star\star$]\label{exo:pl:trade-capture-and-booking:5}
The confirmation system is wrong about 60 trades. Which events did it miss, and what would it have had to subscribe to?
\end{exercise}

\begin{exercise}[$\star\star$]\label{exo:pl:trade-capture-and-booking:6}
A trade is booked in the wrong book and moved the next day. How should the move be recorded, and why?
\end{exercise}

\begin{exercise}[$\star\star\star$]\label{exo:pl:trade-capture-and-booking:7}
\emph{Coding.} Rerun \texttt{downstream} with the risk system copying at 12:00 as well as at 18:00. How many trades is it wrong about?
\end{exercise}

\begin{exercise}[$\star\star\star$]\label{exo:pl:trade-capture-and-booking:8}
\emph{Find the flaw.} ``Every system has the trade, so they must agree: each one reads it from the trading system when it needs it.''
\end{exercise}

\section{Problem: Three Versions on Friday}

\begin{problem}[{Weekend problem --- copies against events}]\label{pb:pl:trade-capture-and-booking:1}
The block, the 200 swaps and the week of the chapter.

\medskip\noindent\textbf{Part I --- The model.}
\begin{enumerate}
\item What does the canonical model hold, and why the instrument as the pricing library serialises it?
\item How is a UTI built?
\item What does the validator check?
\item What does a \omterm{def:pl:trade-capture-and-booking:model}{booking model} decide?
\item Why is a trade never edited?
\end{enumerate}

\medskip\noindent\textbf{Part II --- The week.}
\begin{enumerate}[resume]
\item What events happened to the swaps?
\item What did the hook's swap go through, and what did each system hold on Friday?
\item What does each copy policy miss?
\item How many trades was each system wrong about?
\item On how many trades did the systems disagree, and how many under events?
\end{enumerate}

\medskip\noindent\textbf{Part III --- The block.}
\begin{enumerate}[resume]
\item How was the block executed, and at what average price?
\item What are the exact shares of the three funds?
\item What does each rule allocate, and with what residual?
\item What is the largest deviation from target under each rule?
\item What price does each fund pay?
\end{enumerate}

\medskip\noindent\textbf{Part IV --- The verdict.}
\begin{enumerate}[resume]
\item State the \emph{named result}: the number of downstream disagreements after a week of lifecycle events under copy-based and event-based
distribution, and the allocation residuals of a block split among three funds under three rounding rules.
\item Which allocation rule would you adopt, and why?
\item What would you require of a new downstream system?
\item What does \omterm{def:pl:trade-capture-and-booking:stp}{straight-through processing} measure?
\item In one sentence: what is a trade?
\end{enumerate}
\end{problem}

\section{Interview questions}

\begin{interviewq}[$\star$ \iqroles{developer}]\label{iq:pl:trade-capture-and-booking:1}
Three systems show three versions of the same trade. How do you find out which is right, and how do you stop it happening?
\end{interviewq}

\begin{interviewq}[$\star\star$ \iqroles{developer}]\label{iq:pl:trade-capture-and-booking:2}
How would you represent a trade so that pricing, risk, confirmation and settlement can all use it?
\end{interviewq}

\begin{interviewq}[$\star\star$ \iqroles{developer}]\label{iq:pl:trade-capture-and-booking:3}
A block is filled at twenty prices for five funds. How do you allocate it?
\end{interviewq}

\begin{interviewq}[$\star\star$ \iqroles{developer}]\label{iq:pl:trade-capture-and-booking:4}
What is a novation, and what does it change in the systems that hold the trade?
\end{interviewq}

\begin{interviewq}[$\star\star\star$ \iqroles{developer}]\label{iq:pl:trade-capture-and-booking:5}
Design the \omterm{def:pl:trade-capture-and-booking:capture}{trade capture} and \omterm{def:pl:trade-capture-and-booking:capture}{booking} layer of a multi-asset trading firm.
\end{interviewq}
